\documentclass[10pt,a4paper,sans]{moderncv}
\moderncvstyle{classic}
\moderncvcolor{blue}
\usepackage[utf8]{inputenc}
\usepackage[scale=0.85]{geometry}
\usepackage{helvet}

\name{Alexandro}{Disla}
\title{Suivi, Évaluation et Apprentissage --- Information Stratégique en Santé}
\address{Rue Saint-Louis Jeanty \#14, Route de Frère}{Port-au-Prince, Haïti}{}
\phone[mobile]{(+509) 4148-3700}
\email{alexandrodisla@hotmail.com}
\extrainfo{GitHub: \url{https://github.com/AD0791}}

\begin{document}
\makecvtitle

\section{Profil Professionnel}
\cvitem{}{Plus de dix ans en suivi, évaluation et systèmes d'information, dont \textbf{trois ans comme Officier de Suivi \& Évaluation sur un programme VIH} à la Caris Foundation International --- rapportage de routine de cliniques multi-sites et remontée des \textbf{indicateurs MER à l'USAID via DATIM}, la plateforme du PEPFAR bâtie sur DHIS2. Mon travail porte sur la part du suivi-évaluation qui doit résister à la vérification : protocoles d'assurance qualité des données établis et supervisés, chiffres remontés jusqu'au registre source, collecte numérique standardisée pour que la donnée soit juste à la saisie plutôt que réparée ensuite. Économiste appliqué de formation et \textbf{ingénieur de données} en exercice, je construis et administre les bases, les pipelines et les tableaux de bord dont dépend un système d'information sanitaire --- la réconciliation entre systèmes qui ne se parlent pas est un travail que j'exerce professionnellement. Résidant à Port-au-Prince ; français et créole haïtien maternels, anglais courant.}

\section{Compétences Clés}
\cvitem{}{
\begin{itemize}
    \item \textbf{Cadres MEL \& suivi de performance :} Cadres logiques, chaînes de résultats, plans de suivi et tableaux de suivi des indicateurs ; définition d'indicateurs avec valeurs de référence, cibles, désagrégations et sources de vérification ; cycles de rapportage mensuel, trimestriel, semestriel et annuel depuis des institutions sanitaires multi-sites.
    \item \textbf{Qualité des données \& préparation aux audits :} Protocoles DQA établis et supervisés, vérification par remontée aux registres sources, suivi des actions correctives jusqu'à clôture, dédoublonnage des listes de bénéficiaires --- la discipline exacte qu'exige une vérification indépendante.
    \item \textbf{Systèmes d'information \& réconciliation de données :} Administration de bases relationnelles en accès concurrent, pipelines ETL, contrôle d'accès par rôle et stockage sécurisé ; rapprochement de sources hétérogènes et résolution d'identités entre systèmes ; rapportage bailleur via DATIM/DHIS2 en tant que producteur de données.
    \item \textbf{Analyse \& visualisation :} Power BI, Looker Studio, Quarto et mWater pour un suivi périodique des indicateurs ; SQL et Python pour l'automatisation du rapportage ; modélisation statistique et échantillonnage.
    \item \textbf{Renforcement de capacités \& coordination :} Supervision d'agents de collecte multi-sites, encadrement technique de personnel de suivi-évaluation, conception et animation de formations aux outils et à la qualité des données ; coordination avec le MSPP, la DINEPA et le MEF.
\end{itemize}}

\section{Expériences en Suivi-Évaluation Sanitaire et Systèmes d'Information}

\cventry{Oct 2021 -- Août 2024}{Officier de Suivi \& Évaluation --- programme VIH}{Impact Youth Project, Caris Foundation International}{Haïti}{}{
\begin{itemize}
    \item \textbf{Rapportage sanitaire de routine :} Conduite du cycle de suivi d'un programme VIH sur des cliniques multi-sites, avec consolidation des données des sites dans les délais et remontée des \textbf{indicateurs MER à l'USAID via DATIM}, la plateforme de rapportage du PEPFAR bâtie sur DHIS2.
    \item \textbf{Assurance qualité (DQA) :} Établissement et supervision des protocoles de Data Quality Assessment sur les données intégrées MySQL --- remontée des chiffres rapportés jusqu'aux registres sources et suivi des actions correctives jusqu'à clôture plutôt que jusqu'à leur simple consignation.
    \item \textbf{Systèmes de collecte :} Mise en place et administration des systèmes de collecte numérique alimentant ce rapportage (CommCare, HIVHaiti, MySQL), avec standardisation des formulaires et des flux de données entre sites.
    \item \textbf{Encadrement :} Supervision des agents de collecte sur l'ensemble des sites et encadrement technique du personnel de suivi-évaluation sur les procédures de collecte et les contrôles qualité.
    \item \textbf{Automatisation \& tableaux de bord :} Rapports d'activité développés en Python et SQL, réduisant le temps de traitement manuel de 40 \% ; tableaux de bord programme maintenus sur AWS (Quarto) et Power BI pour le suivi périodique des indicateurs.
\end{itemize}}

\cventry{Jan 2026 -- Mars 2026}{Analyste de Données \& Consultant MEAL}{Anseye Pou Ayiti}{Haïti}{}{
\begin{itemize}
    \item Définition de la stratégie de collecte et mise en œuvre des indicateurs de performance de trois programmes, de la définition de l'indicateur jusqu'à sa restitution.
    \item Codage des instruments XLSForm pour ODK avec logiques de saut et contraintes de validation, administration du système ODK, et maintenance des scripts d'intégration Python sur AWS.
\end{itemize}}

\cventry{Oct 2023 -- Jan 2026}{Consultant mWater \& Analyste de Données}{HANWASH}{Haïti / Télétravail}{}{
\begin{itemize}
    \item Raffinement des plans MEAL et des tableaux de suivi des indicateurs alignés sur les normes JMP (OMS/UNICEF) pour des programmes d'eau, d'hygiène et d'assainissement.
    \item Déploiement d'enquêtes géospatiales sur mWater, administration de la base, et tableaux de bord automatisés (mWater, Power BI) connectés aux API de collecte ; rédaction des guides techniques et formation des équipes.
\end{itemize}}

\cventry{Nov 2015 -- Août 2023}{Économiste --- Direction de Planification Économique et Sociale (DPES)}{Ministère de la Planification et de la Coopération Externe (MPCE)}{Port-au-Prince}{}{Contribution aux rapports de suivi de l'exécution du Plan Triennal d'Investissement 2014--2016, avec suivi d'indicateurs de projets publics sur l'ensemble des secteurs et des départements ; affectation à l'unité statistique et modélisation et participation à la construction du << Modèle de Planification du Développement >>, outil de prévision, de simulation et d'analyse d'impact des politiques publiques ; suivi des indicateurs macroéconomiques avec le Ministère de l'Économie et des Finances.}

\section{Expériences en Administration de Bases de Données et Systèmes}

\cventry{Févr 2026 -- Présent}{Ingénieur Logiciel (Fullstack)}{Tekkod LLC}{Télétravail}{}{Direction d'un programme de modernisation d'infrastructure mobile, avec fiabilité hors ligne et sécurité des données pour un usage terrain ; tableaux de bord Looker Studio avec recherche et filtrage avancés pour la visibilité opérationnelle.}

\cventry{Mars 2021 -- Mai 2024}{Développeur Backend \& Ingénieur de Données}{Tekkod LLC}{Télétravail}{}{Administration de bases relationnelles en accès concurrent et résolution des défaillances qui rendent un système inutilisable sur le terrain ; pipelines ETL (Apache Beam) migrant les données MongoDB vers BigQuery ; services REST sécurisés (FastAPI, JWT) avec contrôle d'accès par rôle et documentation OpenAPI.}

\cventry{Nov 2015 -- Mars 2021}{Consultant Logiciel \& Spécialiste Intégration de Données}{CassionSoft (Projet UNOPS)}{Haïti}{}{Intégration de données sur un projet sous contrat des Nations Unies, avec exposition directe aux procédures et aux standards documentaires onusiens ; authentification RBAC (OAuth2/JWT) et validation de la logique d'intégration par tests unitaires.}

\section{Formation}
\cvitem{2010--2014}{Études Supérieures en Économie Appliquée --- C.T.P.E.A (Centre de Techniques de Planification et d'Économie Appliquée), Port-au-Prince : cycle de quatre ans en économie et statistique appliquées. \textit{Scolarité complétée ; Diplôme d'Études Supérieures en attente de la soutenance du mémoire de sortie (attestation officielle disponible).}}
\cvitem{2009--2010}{Baccalauréat (Mention) --- Institution Saint-Louis de Gonzague, Port-au-Prince}
\cvitem{Autoformation}{FlatIron School (Software Engineering), CodeWithMosh, Udemy}

\section{Compétences Techniques}
\cvitem{Info. sanitaire}{\textbf{DHIS2 / DATIM (rapportage des indicateurs MER)}, CommCare, HIVHaiti, mWater, protocoles DQA, tableaux de suivi des indicateurs}
\cvitem{Collecte mobile}{\textbf{ODK (Open Data Kit), KoboToolbox, XLSForm}, logiques de saut et contraintes de validation, conception et gestion d'enquêtes}
\cvitem{Analyse \& reporting}{\textbf{Excel avancé}, \textbf{Power BI (expert)}, SQL, Python (Pandas), R (Quarto), Looker Studio, modélisation statistique et échantillonnage}
\cvitem{Infrastructure}{MySQL, PostgreSQL, MongoDB, SQLite, Azure, AWS, Google Cloud (BigQuery), Docker, Git}
\cvitem{Langues}{\textbf{Français (maternel)}, \textbf{Créole haïtien (maternel)}, \textbf{Anglais (courant)}}

\section{Référence Professionnelle}
\cvitem{}{Davidson Adrien --- ancien Directeur Suivi \& Évaluation, Caris Foundation International --- (+509) 4460-6638}

\end{document}
